\section{Ten rank-ten exclusions from central reciprocity}
The rings
\[
R10\_1\_k,\qquad
k\in\{65,89,90,91,120,143,144,147,162,182\},
\]
have no categorification over $\mathbb C$.
The separately verified signed/Galois reduction in \nolinkurl{work/general_signed}
shows that a hypothetical complex categorification of each listed based
ring would give a pseudo-unitary categorification of the same based ring.
The following obstruction rules out that consequence.

Let $\mathcal C$ be pseudo-unitary, $D=\operatorname{FPdim}(\mathcal C)$,
$F:Z(\mathcal C)\to\mathcal C$ the forgetful functor, and $I$ its adjoint.
For each one-dimensional representation $\chi$ of its complexified fusion
ring, let $f_\chi=\sum_i\chi(X_i)\chi(X_i^*)$ be its formal codegree.
There is a distinct simple $A_\chi\subset I(1)$ with
\[
[F(A_\chi):1]=1,\qquad
\dim A_\chi=\beta_\chi=D/f_\chi.
\]
Moreover, for every central simple $Z$,
\[
S_{A_\chi,Z}=\beta_\chi\chi(F(Z)).
\]
These are Ostrik's induction/formal-codegree theorem and the identification
of the corresponding central character (Ostrik, \emph{Pivotal fusion categories of rank 3},
arXiv:1309.4822, Example 2.9 and the discussion preceding
Theorem 2.13; equivalently Deshpande, arXiv:1611.04486v1, equations
(35)--(36)). Symmetry of $S$ gives the necessary equations
\begin{equation}\label{eq:r10reciprocity}
\beta_\chi\chi(v_\psi)=\beta_\psi\psi(v_\chi),
\qquad v_\chi=[F(A_\chi)].
\end{equation}
Only scalar characters are used, also for the eight noncommutative rings;
no noncommutative character recovery or normalization is needed.

Write $N_{ij}^k$ for the supplied based tensor and define its exact
induction Gram matrix by
\[
B_{ij}=\sum_{a,k}N_{aj}^kN_{ka^*}^{i}.
\]
Adjunction gives $B=MM^{\mathsf T}$, where the columns of $M$ are the
restriction vectors of all simples of $Z(\mathcal C)$. Hence any distinct
chosen $v_\chi$ satisfy
\[
 v_\chi\in\mathbb Z_{\ge0}^{10},\quad
(v_\chi)_0=1,\quad d\cdot v_\chi=\beta_\chi,
\quad 0\le(v_\chi)_i\le\min(B_{0i},\lfloor\sqrt{B_{ii}}\rfloor),
\]
and $Q=B-\sum_\chi v_\chi v_\chi^{\mathsf T}$ is entrywise nonnegative
and positive semidefinite. Also $\sum_\chi v_\chi\le B_{0,*}$
coordinatewise, since every selected column has zeroth coordinate one.
The verifier uses only the entrywise and $2\times2$ principal-minor
consequences of positivity. These weak consequences already suffice.

For every listed ring, the dimension vector consists of four $1$'s,
two $2+2\sqrt2$'s and four $3+2\sqrt2$'s, and
$D=96+64\sqrt2$. All scalar characters and codegrees lie in
$\mathbb Q(\sqrt2,i)$. The certificate stores them using the primitive
element $t=\sqrt2+i$, whose minimal polynomial is $t^4-2t^2+9$.
The independent verifier decodes them into the rational basis
$(1,\sqrt2,i,i\sqrt2)$ and checks each of the $100$ multiplicativity
equations for each distinct character. It checks $f_\chi\beta_\chi=D$
exactly. It independently regenerates every possible $v_\chi$ from the
two rational dimension equations and the displayed finite bounds.
It then exhausts equations~\eqref{eq:r10reciprocity}, the Gram bounds,
and first-row inequalities. All ten finite searches are empty, after
only three or five search nodes each. It does not assume that the stored
scalar table is complete, nor does it impose first-row equality.

The machine receipt is \nolinkurl{exclusion_receipts.json}. The initial
searches are \nolinkurl{reciprocal_partial.py} for the eight noncommutative
rings and \nolinkurl{reciprocal.py} for the two commutative rings 89 and 147.
The separate, standard-library-only replay is
\nolinkurl{work/audit_new_models/verify_rank10_reciprocity.py}; its complete
output is \nolinkurl{work/audit_new_models/rank10_reciprocity_audit.json}.
The signed reduction was replayed separately using the verifier in
\nolinkurl{work/general_signed}; it is not part of this standard-library
scalar-character verifier. Thus the
exclusions concern arbitrary complex fusion categories, not just
pseudo-unitary or spherical categories assumed in advance.
